\begin{abstractEng}{eBPF; open-source AI inference framework; cross-language program analysis; call chain tracing; Ollama; HybriChain}
Open-source large language models are now widely deployed in enterprise private environments. Inference engines like Ollama use a Go-CGO-C++ tri-layer hybrid architecture with dual-process decoupling. This design creates significant observability blind spots at the system level. When attackers exploit API endpoints for directory traversal or inject malformed payloads, security operators have no way to trace how external data flows through the inference engine internally, nor can they identify which kernel-level system calls the behavior triggers.

This thesis presents HybriChain---a cross-language call chain tracing and security analysis system built for multi-language inference engines. The cross-language symbol mapping problem is addressed through a five-layer symbol classification system (f$_0$--f$_4$), which unifies function symbols from different namespaces---Go runtime, CGO ABI, and low-level computing libraries---into a shared semantic space. Tracing overhead is minimized by an eBPF-based probe sinking strategy that attaches uprobes only to core inference functions. Noisy data is handled by a five-step funnel noise reduction pipeline compressing raw event streams by 2--3 orders of magnitude. For call chain stitching, a timestamp-aligned confidence algorithm assigns Confirmed, Inferred, or Unknown states to each cross-language edge. On top of these components, an anomaly behavior analysis engine detects high-risk behaviors such as futex blocking, abnormal mmap allocation, and sensitive path access.

Experiments using Ollama as the target validate the system. HybriChain reconstructs complete cross-language call chains from raw kernel event streams and identifies anomalous behaviors in inference engines. The results suggest a practical approach for observability and security auditing of open-source AI inference software.
\end{abstractEng}
